\section{Abstract}

Synthetic minimal cells define the gene set required for life, yet how their genomes are expressed at the RNA level remains largely uncharted.
We combined full-length long-read RNA sequencing, short-read quantification, and matched proteomics to map transcription, RNA processing, and translation across \jsyn~(\syn) and its minimized derivative \jsyna~(\syna).
Long reads resolved 459 operons in \syn~with matched promoters and terminators, whose transcripts were pervasively reshaped by processing biased toward the 3$'$ end.
Transcript level was the dominant determinant of protein abundance, and the antisense and intergenic transcription that complicated \syn~traced to inherited mis-annotation and synthetic artifacts that genome minimization preferentially removed.
Halving the genome preserved gene order yet functionally reallocated the transcriptome toward the translation machinery, while RNA polymerase and central-metabolic enzymes declined, coherent with the slower growth of \syna.
These data provide the RNA-level foundation for whole-cell modeling of the minimal cell.

\newpage
